\section{Begründung zentraler Entwicklungsentscheidungen}

Im Verlauf der Entwicklung wurden mehrere grundlegende Entscheidungen getroffen, die den Charakter, die Architektur und die technische Umsetzung von Pablo maßgeblich geprägt haben.
Im Folgenden werden diese Entscheidungen jeweils mit den betrachteten Alternativen und der jeweiligen Begründung dargestellt.

\subsection{Physischer Companion statt reiner Softwareassistent}

Eine erste zentrale Entscheidung betraf die Frage, ob Pablo als reine Softwarelösung oder als physisches Produkt umgesetzt werden soll.
Softwarebasierte Assistenten lassen sich mit geringerem Aufwand entwickeln und auf bestehenden Endgeräten betreiben.
Demgegenüber bietet ein physisches System eine deutlich höhere Präsenz im Alltag der Nutzenden.
Forschungsarbeiten im Bereich der sozialen Robotik belegen, dass verkörperte Systeme eine stärkere emotionale Bindung erzeugen und als vertrauenswürdiger wahrgenommen werden als rein virtuelle Agenten.
Die Entscheidung fiel daher bewusst zugunsten eines physischen Companion-Bots, um Pablo einen eigenständigen Produktcharakter zu verleihen und eine persönliche Bindung zu ermöglichen, die über die reine Funktionalität hinausgeht.

\subsection{Modular erweiterbare Plugin-Struktur}

Für die Softwarearchitektur wurde zwischen einer monolithischen Implementierung mit fest integrierten Funktionen und einer modularen Plugin-Struktur mit dynamischer Erweiterbarkeit abgewogen.
Eine monolithische Lösung wäre in der initialen Entwicklung einfacher umsetzbar gewesen, hätte jedoch jede funktionale Erweiterung an eine Veränderung des Kernprozesses gebunden.
Die Entscheidung fiel auf eine modulare API-basierte Struktur, da diese es ermöglicht, neue Funktionen nachträglich zu ergänzen, ohne bestehenden Code verändern zu müssen.
Darüber hinaus erhöht die modulare Architektur die Anpassbarkeit an unterschiedliche Zielgruppen und Nutzungsszenarien und schafft eine Grundlage für ein langfristiges Open-Source-Ökosystem.

\subsection{Raspberry Pi als zentrale Recheneinheit}

Für die zentrale Recheneinheit wurden verschiedene Plattformen in Betracht gezogen, darunter leistungsfähigere Einplatinencomputer sowie Cloud-basierte Ansätze, bei denen die gesamte Verarbeitung auf einem externen Server erfolgt.
Die Entscheidung fiel auf einen Raspberry Pi, da dieser eine gute Verfügbarkeit, eine breite Community-Unterstützung und eine ausreichende Schnittstellenvielfalt für die Anbindung von Mikrofon, Display, Lautsprecher und Mikrocontrollern bietet.
Obwohl die Rechenleistung im Vergleich zu leistungsstärkeren Alternativen begrenzt ist, reicht sie für die lokale Koordination der einzelnen Systemkomponenten aus.
Zudem eignet sich der Raspberry Pi aufgrund seiner kompakten Bauform besonders gut für den Einsatz in einem Prototyp mit begrenztem Gehäusevolumen.

\subsection{Wakeword- und Sprachsteuerung}

Für die Aktivierung und Steuerung des Systems wurden neben der Sprachsteuerung auch physische Eingabemethoden wie Tasten oder Touchscreen-Bedienung erwogen.
Die Entscheidung fiel auf eine wakeword-basierte Sprachsteuerung, da diese eine niedrigschwellige, freihändige Interaktion ermöglicht.
Nutzende können Pablo aktivieren und Befehle erteilen, ohne ihre aktuelle Tätigkeit unterbrechen zu müssen.
Dieser Ansatz unterstützt eine möglichst natürliche Zugangsmöglichkeit zum System und entspricht dem Companion-Konzept, bei dem die Interaktion intuitiv und beiläufig erfolgen soll.

\subsection{Visuelle Persona}

Bei der Gestaltung der Benutzeroberfläche wurde zwischen einer rein funktionalen Darstellung, etwa einer textuellen Statusanzeige, und einer visuellen Persona in Form eines animierten Charakters abgewogen.
Die Entscheidung fiel auf die Darstellung eines Pixel-Art-Faultiers als visuelle Identität von Pablo.
Diese Gestaltung erzeugt einen hohen Wiedererkennungswert, ermöglicht eine freundlichere und zugänglichere Interaktion und grenzt Pablo bewusst von rein technischen Assistenzsystemen ab.
Durch zustandsabhängige Animationen kann zudem der aktuelle Systemstatus auf eine intuitive und ansprechende Weise kommuniziert werden.

\subsection{Kompaktes integriertes Display}

Alternativ zur Integration eines eigenen Displays wäre eine Ausgabe ausschließlich über Sprache oder über ein externes Gerät wie ein Smartphone möglich gewesen.
Die Entscheidung für ein kompaktes, direkt im Gehäuse verbautes Display wurde getroffen, um Pablo ein lokales visuelles Feedback zu ermöglichen.
Das Display dient der Darstellung des Avatars, der Anzeige von Statusinformationen und der visuellen Begleitung der Sprachinteraktion.
Durch die Integration des Displays in das Gehäuse wird Pablo als eigenständiges, in sich geschlossenes Produkt wahrnehmbar, das keine zusätzlichen externen Geräte für die Nutzung benötigt.